\section{Mô hình hóa đồ thị và quan hệ lân cận}

Để xử lý bằng thuật toán, bàn Minesweeper kích thước $R \times C$ được mô hình hóa thành một đồ thị vô hướng $G = (V, E)$.

\textbf{Tập đỉnh (V):}  
Mỗi ô trên bàn là một đỉnh, với tọa độ $(r, c)$, trong đó $0 \le r < R$, $0 \le c < C$.  
Tổng số đỉnh là $R \times C$.

\textbf{Tập cạnh (E):}  
Hai đỉnh được nối với nhau nếu chúng kề nhau theo 8 hướng (trên, dưới, trái, phải và 4 đường chéo). Điều này tương ứng với lân cận Moore trong đồ thị lưới.

Vì mỗi ô chỉ có tối đa 8 ô xung quanh, nên bậc của đồ thị bị chặn trên bởi một hằng số (tối đa 8). Điều này rất quan trọng khi phân tích độ phức tạp.

Mỗi đỉnh trong mô hình sẽ lưu các trạng thái rời rạc như: có mìn hay không, đã mở hay chưa, có cắm cờ hay không, và số mìn lân cận.

